\documentclass[11pt,letterpaper]{article}
\def\SuiteDocTitle{Role Modeling and Story Discovery — Workbook Template}
\def\SuiteDocCode{WB-ROLE-A}
% Shared professional style: pdfLaTeX compatible, no external font files.
\usepackage[T1]{fontenc}
\usepackage[utf8]{inputenc}
\usepackage{lmodern}
\usepackage{microtype}
\usepackage[margin=0.83in,headheight=16pt,footskip=28pt]{geometry}
\usepackage[table]{xcolor}
\definecolor{navy}{HTML}{12304A}
\definecolor{teal}{HTML}{087F83}
\definecolor{gold}{HTML}{B78633}
\definecolor{blueback}{HTML}{F0F5F8}
\definecolor{tealback}{HTML}{EEF8F6}
\definecolor{amberback}{HTML}{FFF7E8}
\definecolor{grayback}{HTML}{F6F7F8}
\usepackage{array,tabularx,longtable,booktabs}
\usepackage{enumitem}
\usepackage{amssymb}
\usepackage{titlesec}
\usepackage{fancyhdr,lastpage}
\usepackage{tcolorbox}
\tcbuselibrary{breakable}
\usepackage{pgffor}
\usepackage{hyperref}
\hypersetup{colorlinks=true,linkcolor=teal,urlcolor=teal,pdfauthor={Engineering Learning Series},pdftitle={\SuiteDocTitle}}
\setlength{\parindent}{0pt}
\setlength{\parskip}{5.5pt}
\setlist{topsep=3pt}
\setcounter{tocdepth}{2}
\titleformat{\section}{\Large\bfseries\color{navy}}{\thesection}{0.7em}{}
\titleformat{\subsection}{\large\bfseries\color{teal}}{\thesubsection}{0.6em}{}
\titleformat{\subsubsection}{\normalsize\bfseries\color{navy}}{\thesubsubsection}{0.6em}{}
\pagestyle{fancy}
\fancyhf{}
\fancyhead[L]{\scriptsize\color{navy}AGILE REQUIREMENTS ENGINEERING}
\fancyhead[R]{\scriptsize\color{teal}\SuiteDocCode}
\fancyfoot[L]{\scriptsize Independent learning and implementation materials}
\fancyfoot[R]{\scriptsize Page \thepage\ of \pageref{LastPage}}
\renewcommand{\headrulewidth}{0.35pt}
\renewcommand{\footrulewidth}{0pt}
\newcommand{\SuiteHeader}[2]{%
{\color{teal}\small\bfseries ENGINEERING DOCUMENTATION SUITE \hfill #2}\par\vspace{8pt}
{\color{navy}\LARGE\bfseries\SuiteDocTitle}\par\vspace{5pt}
{\color{black!68}\normalsize #1}\par\vspace{8pt}
{\color{gold}\rule{\linewidth}{2pt}}\par\vspace{8pt}}
\newcommand{\WritingLines}[1]{\par\foreach \n in {1,...,#1}{\noindent\rule{\linewidth}{0.35pt}\par\vspace{0.28cm}}}
\newcommand{\FreeSpace}[1]{\begin{tcolorbox}[colback=white,colframe=black!25,height=#1,boxrule=0.5pt]\end{tcolorbox}}
\newcommand{\Field}[1]{\noindent\textbf{#1}\quad\rule{0.42\linewidth}{0.35pt}\par\vspace{0.3cm}}
\newcommand{\CheckItem}[1]{\noindent$\square$\quad #1\par\vspace{2pt}}
\newcommand{\SmallNote}[1]{{\footnotesize\color{black!62}#1}\par}
\newcommand{\Guidance}[1]{\begin{tcolorbox}[colback=blueback,colframe=teal,breakable,title=Guidance]#1\end{tcolorbox}}
\sloppy

\begin{document}
\SuiteHeader{Facilitator instructions | examples and working space}{Workbook template}
\tableofcontents\medskip\noindent\rule{\linewidth}{0.4pt}\medskip
\section{Workbook objective}
Facilitate a repeatable workshop with shared decisions and tangible evidence. This annotated edition contains original illustrative answers. Use the paired blank edition for actual fieldwork and printing.

\section{Participants and inputs}
\begin{itemize}[leftmargin=*,itemsep=2pt]
\item Product decision-maker and representative users
\item Engineers and test specialists
\item Existing observations, goals and backlog items
\item Time-boxed session, typically 45–90 minutes per working segment
\end{itemize}
\section{Step 1: Scope the opportunity}
Name the decision and the user outcome to investigate.

\begin{tcolorbox}[colback=tealback,colframe=navy,title={Illustration},breakable]
A requester should learn whether a submitted item is waiting on their action.

\end{tcolorbox}
\subsection{Record your analysis}
\WritingLines{4}
\section{Step 2: Inventory candidate roles}
Brainstorm distinct behavior patterns and record source of evidence.

\begin{tcolorbox}[colback=tealback,colframe=navy,title={Illustration},breakable]
Requester, reviewer, service coordinator, auditor.

\end{tcolorbox}
\subsection{Record your analysis}
\WritingLines{4}
\section{Step 3: Consolidate and differentiate}
Combine equivalent roles but preserve meaningful differences.

\begin{tcolorbox}[colback=tealback,colframe=navy,title={Illustration},breakable]
Auditor is read-only while reviewer can change status.

\end{tcolorbox}
\subsection{Record your analysis}
\WritingLines{4}
\section{Step 4: Design an elicitation protocol}
Use observation and neutral interview questions.

\begin{tcolorbox}[colback=tealback,colframe=navy,title={Illustration},breakable]
Observe one delayed request and ask what the requester expected next.

\end{tcolorbox}
\subsection{Record your analysis}
\WritingLines{4}
\section{Step 5: Capture findings by evidence tier}
Separate observation, claim, decision and assumption.

\begin{tcolorbox}[colback=tealback,colframe=navy,title={Illustration},breakable]
Observed: people ask support for status; assumed: reminders reduce requests.

\end{tcolorbox}
\subsection{Record your analysis}
\WritingLines{4}
\section{Step 6: Create goal-to-story map}
Connect each role to goals and candidate slices.

\begin{tcolorbox}[colback=tealback,colframe=navy,title={Illustration},breakable]
Requester -> track status -> view current state and history.

\end{tcolorbox}
\subsection{Record your analysis}
\WritingLines{4}
\section{Step 7: Validate with stakeholders}
Review findings with actual users and accountable product representative.

\begin{tcolorbox}[colback=tealback,colframe=navy,title={Illustration},breakable]
Ask two requesters to classify the proposed status states.

\end{tcolorbox}
\subsection{Record your analysis}
\WritingLines{4}
\section{Workshop close-out}
\begin{itemize}[leftmargin=*,itemsep=2pt]
\item Summarize decisions, unresolved assumptions and owners.
\item Create or update linked records.
\item Agree the verification event before the next commitment.
\end{itemize}

\end{document}
